\documentclass[10pt,letterpaper]{article}
\usepackage[optengjnl]{opteng}
\usepackage{booktabs}

\subjterm{optical sensing, calibration, field instrumentation}
\loginfo{Example manuscript compiled \today}
\doi{10.0000/example.opteng}

\begin{document}

\title{Compact optical sensing for resilient field measurements}
\author{Avery Morgan}
\address{Centre for Applied Photonics, Northbridge Institute}
\email{avery.morgan@example.org}

\begin{abstract}
Field instruments must remain useful when temperature, vibration, and available
power change together. We describe a compact optical sensing workflow that uses
an internal reference path and a short calibration sequence. In a laboratory
trial the reference reduced the median measurement drift by 38 percent while
adding less than two seconds to each acquisition. The example is intentionally
small, so the complete manuscript can be adapted to the page-length checks used
by Optical Engineering.
\end{abstract}

\section{Introduction}
Optical instruments can measure a large range of physical quantities without
placing electrical contacts in the sensing region. Their performance in the
field nevertheless depends on alignment, ambient temperature, and the time
between calibration points. A useful design therefore needs both a sensitive
measurement channel and a modest way to estimate its current error.

\section{Method}
The instrument alternates a sample measurement with a reference measurement.
Let $s$ be the sample response, $r$ the reference response, and $s_0,r_0$ the
values recorded during calibration. We report the normalized response
\begin{equation}
  q = \frac{s/r}{s_0/r_0},
  \label{eq:normalized-response}
\end{equation}
which removes a common multiplicative change in source intensity.

\begin{figure}[t]
  \centering
  \fbox{\rule{0pt}{28mm}\rule{0.82\columnwidth}{0pt}}
  \caption{Editable placeholder for the optical path diagram. Replace the
  framed area with a publication-ready figure before submission.}
  \label{fig:path}
\end{figure}

\section{Results}
Table~\ref{tab:trial} shows the synthetic pilot measurements used to exercise
the template's compact table style.

\begin{table}[t]
  \centering
  \caption{Illustrative calibration trial.}
  \label{tab:trial}
  \begin{tabular}{lrr}
    \toprule
    Condition & Drift (\%) & Time (s)\\
    \midrule
    Uncorrected & 6.4 & 1.8\\
    Reference path & 4.0 & 3.5\\
    \bottomrule
  \end{tabular}
\end{table}

The reference path improves stability at the cost of one additional acquisition.
This trade-off is attractive when a later recalibration is expensive or when
the instrument cannot be returned to a controlled laboratory.

\section{Discussion and conclusion}
Equation~\ref{eq:normalized-response} is a simple model, not a substitute for
the uncertainty analysis required by a real experiment. The starter provides
the structural pieces needed to add that analysis: a clear method section,
editable visual material, a compact results table, and a bibliography. Authors
should replace every synthetic value and placeholder metadata before submission.

\begin{thebibliography}{9}
\bibitem{saleh}
B. E. A. Saleh and M. C. Teich, \emph{Fundamentals of Photonics}, 3rd ed.,
Wiley, 2019.

\bibitem{holst}
G. C. Holst, \emph{Electro-Optical Imaging System Performance}, JCD Publishing,
2008.
\end{thebibliography}

\end{document}
